\section{参数影响与结果讨论}

\subsection{人工粘性的影响}

对比关闭和开启人工粘性的结果，讨论：

\begin{itemize}[leftmargin=2em]
    \item 间断附近振荡是否减弱；
    \item 激波是否被过度抹宽；
    \item 接触间断的分辨率变化；
    \item 不同人工粘性系数的影响。
\end{itemize}

\figureplaceholder{不同人工粘性系数的结果比较}{fig:viscosity}

\subsection{CFL 的影响}

建议至少比较三组满足稳定条件的 CFL，例如 0.2、0.5 和 0.8。

\begin{table}[H]
    \centering
    \caption{不同 CFL 的计算表现}
    \label{tab:cfl}
    \begin{tabular}{cccc}
        \toprule
        CFL & 时间步数 & 是否稳定 & 主要现象\\
        \midrule
        0.2 & -- & -- & --\\
        0.5 & -- & -- & --\\
        0.8 & -- & -- & --\\
        \bottomrule
    \end{tabular}
\end{table}

\subsection{网格分辨率的影响}

比较不同 $N_x$，讨论激波厚度、误差和计算成本。

\figureplaceholder{不同网格分辨率的结果比较}{fig:grid}

\subsection{数值异常与适用范围}

记录程序运行中观察到的负压、负密度、NaN、步数上限或边界影响。
区分以下两类结论：

\begin{itemize}[leftmargin=2em]
    \item 格式本身的限制；
    \item 当前程序实现仍需改进的部分。
\end{itemize}
